\section{Experiments}
\label{sec:experiments}

\subsection{Experimental Setup}

\paragraph{Model and Predictor.}
We evaluate on LLaMA-3.1-8B~\cite{grattafiori2024llama3}, which uses SwiGLU FFN blocks~\cite{shazeer2020glu} with $d_{\text{model}}=4096$ and $d_{\text{ff}}=14336$ across 32 layers.
The sparsity predictor is a two-layer MLP with 2.37M parameters per layer (75.9M total), implemented as Linear($4096 \to 128$) + GELU + Linear($128 \to 14336$), trained to produce per-neuron importance scores that minimize output-distribution divergence under sparsity.

\paragraph{Training.}
All predictors are trained for 1000 steps on WikiText-103 with sequence length 2048 and batch size 4 on a single NVIDIA RTX PRO 6000 Blackwell Server Edition GPU (96\,GB), using AdamW~\cite{loshchilov2019adamw} with learning rate $1 \times 10^{-3}$, weight decay $1 \times 10^{-2}$, and cosine decay.
The KL predictor uses Gumbel-Sigmoid relaxation~\cite{jang2017gumbel} with temperature $\tau=1.0$ annealed to $0.1$ over training.

\paragraph{Evaluation.}
We report perplexity on WikiText-2~\cite{merity2017pointer} (seq\_len=2048) and downstream accuracy on ARC-Challenge (25-shot), WinoGrande (5-shot), Massive Multitask Language Understanding (MMLU; 5-shot), Grade School Math 8K (GSM8K; 8-shot), HellaSwag (10-shot), and Physical Intuition QA (PIQA; 0-shot).

\paragraph{Inference Protocol.}
At inference, the predictor produces per-neuron scores that are passed to TEAL's global top-$k$ allocation mechanism~\cite{liu2025teal}, which distributes the sparsity budget across layers non-uniformly.
We term this \textit{predictor-TEAL}: the predictor replaces TEAL's magnitude-based scoring while retaining its allocation strategy.

\paragraph{Baselines.}
We compare against: (1) \textbf{Vanilla TEAL}~\cite{liu2025teal}, which uses activation magnitudes directly\footnote{TEAL is deterministic; single-run evaluation suffices. KL/BCE predictors require multi-seed evaluation due to random initialization.}; (2) \textbf{BCE predictor} (DejaVu-style~\cite{liu2023dejavu}), trained with per-neuron binary cross-entropy on magnitude-based oracle masks for 1000 and 2700 steps; and (3) \textbf{Uniform allocation}, which distributes the sparsity budget equally across layers.


\subsection{Main Results}

Table~\ref{tab:main} presents the primary comparison. Uniform allocation preserves downstream performance best; global allocation is analyzed to understand allocation-quality tradeoffs (\S\ref{sec:scaling}).
At 30\% sparsity, output-distribution KL under uniform allocation achieves PPL 8.81$\pm$0.02 (3-seed), a 14.5\% reduction over TEAL at matched actual sparsity (8.81 vs.\ 10.30, both at 28.5\%; Appendix~\ref{app:ood}).
Uniform allocation yields stable downstream (MMLU 40.7$\pm$1.3\%, 3-seed; Appendix~\ref{app:stability}), while global top-$k$ exhibits high seed variance (35.7$\pm$9.5\%) from late-layer starvation (L31 retention 28.4\%, Appendix~\ref{app:retention}; \S\ref{sec:scaling}).

\paragraph{Task-Type Split.}
KL shows directional improvements on commonsense reasoning (ARC-c, WG, HellaSwag, PIQA: 4/4 wins at 30\%; Appendix~\ref{app:downstream}) but underperforms on knowledge tasks (MMLU, GSM8K: 0/2).
With three seeds, 95\% CIs overlap for individual tasks; we treat PPL ($\sigma{=}0.02$, non-overlapping) as the primary metric.
The KL predictor learns non-uniform per-layer sparsity, protecting middle layers (L14--18 at 14--18\% vs.\ 32\% mean; Appendix~\ref{app:perlayer_sparsity}).
Constraint interventions (Appendix~\ref{app:constraint}) confirm a task-dependent trade-off: removing the learned valley improves MMLU +3.4pp but reduces ARC-c $-$1.5pp, consistent with localized factual associations~\cite{geva2023dissecting}. A C4-trained control (Appendix~\ref{app:c4_control}) does not alleviate the regression, suggesting the deficit is structural to output-distribution optimization rather than training data composition.

At 50\%, the PPL advantage widens (11.92 vs.\ 23.90, 50\% reduction; Figure~\ref{fig:sparsity_sweep}) and generalizes to C4 (Appendix~\ref{app:ood}), but all methods collapse to near-random on MMLU ($\leq$30\%) and GSM8K ($\leq$2.7\%); no allocation strategy recovers at this capacity limit.

\begin{table}[t]
\centering
\caption{Main results on LLaMA-3.1-8B. $^\ast$$2.71\times$ KL compute (\S\ref{sec:isocompute}). $\pm$: 3-seed; unmarked: s42. $^\ddagger$Actual sparsity: KL 28.5\% vs.\ TEAL 29.3\% at 30\% (Appendix~\ref{app:ood}).}
\label{tab:main}
\resizebox{\columnwidth}{!}{
\begin{tabular}{llccccc}
\toprule
Sparsity & Method & PPL$\downarrow$ & ARC-c & WG & MMLU & GSM8K \\
\midrule
0\% & Dense & 6.12 & 58.28 & 78.06 & 65.27 & 48.75 \\
\midrule
\multirow{6}{*}{30\%}
& TEAL uniform & 10.82 & 39.33 & 65.04 & \textbf{45.86} & \textbf{8.79} \\
& BCE global & 11.25 & 41.2 & 66.0 & 36.3 & 6.0 \\
& BCE uniform & 15.09$\pm$0.13 & 33.45 & 61.09 & 43.05 & 3.97 \\
& MSE uniform & 18.52$\pm$1.68 & 35.8$\pm$4.1 & 59.8$\pm$2.6 & 30.7$\pm$2.6 & 2.3$\pm$0.7 \\
& KL global (Ours)$^\ddagger$ & 8.96$\pm$0.02 & \textbf{43.8$\pm$3.6} & \textbf{70.2$\pm$0.9} & 35.7$\pm$9.5 & 8.0$\pm$0.7 \\
& KL uniform$^\ddagger$ & \textbf{8.81$\pm$0.02} & 43.7$\pm$2.5 & 68.9$\pm$0.4 & 40.7$\pm$1.3 & 7.5$\pm$0.1 \\
\midrule
\multirow{5}{*}{50\%}
& TEAL uniform & 23.90 & 27.47 & 54.38 & \textbf{29.26} & 1.74 \\
& BCE (2700 steps)$^\ast$ & 28.13$\pm$0.20 & 27.6$\pm$0.6 & 51.8$\pm$0.6 & 26.1$\pm$0.5 & 2.1$\pm$0.2 \\
& MSE uniform & 229.26 & 25.00 & 50.59 & 23.79 & 1.59 \\
& KL global (Ours) & \textbf{11.84$\pm$0.15} & \textbf{36.3$\pm$3.2} & \textbf{63.8$\pm$0.3} & 28.3$\pm$0.7 & \textbf{2.7$\pm$1.3} \\
& KL uniform & 11.92$\pm$0.07 & 33.8$\pm$1.2 & 62.9$\pm$0.1 & 28.0$\pm$0.6 & 2.2$\pm$0.1 \\
\bottomrule
\end{tabular}}
\end{table}


\subsection{Extended Training Analysis}
\label{sec:isocompute}

At $2.71\times$ KL compute budget (2700 steps, matched wall-clock on a single RTX PRO 6000: $\sim$46 min BCE vs.\ $\sim$45 min KL), BCE achieves PPL 28.13$\pm$0.20 (3-seed), still $2.39\times$ worse than KL (11.79).
The gap is systematic, not due to underfitting: BCE reaches \emph{higher} oracle F1 (0.613 vs.\ 0.540) with 139\% worse PPL, confirming KL optimizes for distributional fidelity rather than magnitude agreement.
This dissociation between oracle agreement and downstream quality is a key finding: better per-neuron classification accuracy does not translate to better output preservation under global allocation.


\subsection{Mask Pattern Analysis}
\label{sec:mask_analysis}

KL-derived masks diverge from magnitude oracles: per-layer Jaccard distance averages 68\% at 50\% (64 samples, 131K tokens; Figure~\ref{fig:mask_heatmap}), with cross-seed variation (36.3\%) confirming systematic divergence (Appendices~\ref{app:cross_seed_jaccard},~\ref{app:mask_viz}).
Despite this, KL achieves PPL 11.79 vs.\ BCE's 28.13$\pm$0.20, confirming that distributional fidelity, not magnitude agreement, produces better masks.
Divergence follows a U-shaped per-layer pattern, peaking in late layers (0.70--0.75 at 50\%).

\begin{figure}[t]
\centering
\includegraphics[width=\columnwidth]{mask_divergence_heatmap.pdf}
\vspace{-8pt}
\caption{Per-layer Jaccard distance between KL-derived and magnitude-based masks at 30\% and 50\% sparsity (64 samples, 131K tokens). Divergence peaks in layers 24--28 ($>$0.73 at 50\%).}
\label{fig:mask_heatmap}
\vspace{-8pt}
\end{figure}

\paragraph{Target-Swap Ablation: Target Quality vs.\ Gradient Coupling.}
To disentangle whether KL's advantage stems from its coupled gradient pathway or from the mask patterns it discovers, we train a standard BCE predictor on KL-derived oracle masks instead of magnitude-based masks---a configuration we label \textbf{C1} (``BCE-on-KL-targets'').
Table~\ref{tab:c1} reports results under TEAL uniform allocation at matched sparsity.

\begin{table}[t]
\centering
\caption{Target-swap ablation (C1): BCE on KL-derived masks vs.\ KL predictor (same seed), TEAL uniform. All seeds show max $\Delta$$\leq$1.70pp at 30\%; max $\Delta$$\leq$0.79pp at 50\% (s42). $^\dagger$\textbf{Provenance:} s42 oracle extracted from an earlier training run (not seed-matched); s123/s456 use seed-matched retrained oracles (Appendix~\ref{app:oracle_provenance}).}
\label{tab:c1}
\resizebox{\columnwidth}{!}{
\begin{tabular}{lllccccc}
\toprule
Sparsity & Seed & Method & PPL$\downarrow$ & ARC-c & WG & MMLU & GSM8K \\
\midrule
\multirow{6}{*}{30\%}
& s42$^\dagger$ & C1 & 9.45 & 40.27 & 67.40 & \textbf{42.30} & \textbf{7.58} \\
& & KL & \textbf{8.82} & \textbf{40.78} & \textbf{68.59} & 40.6 & 7.35 \\
& s123 & C1 & 8.84 & 44.71 & \textbf{69.77} & \textbf{39.92} & 6.82 \\
& & KL & \textbf{8.83} & \textbf{45.31} & 69.38 & 39.4 & \textbf{7.51} \\
& s456 & C1 & 8.82 & 44.71 & \textbf{69.93} & 41.04 & 6.60 \\
& & KL & \textbf{8.79} & \textbf{45.05} & 68.75 & \textbf{42.0} & \textbf{7.58} \\
\midrule
\multirow{2}{*}{50\%}
& s42 & C1 & 11.96 & 32.17 & \textbf{63.14} & 27.79 & \textbf{2.05} \\
& & KL & 11.96 & \textbf{32.59} & 62.35 & \textbf{27.92} & 1.97 \\
\bottomrule
\end{tabular}}
\end{table}

C1 matches KL within 0.79pp at 50\% and $\leq$1.18pp at 30\%\footnote{s42 oracle from earlier checkpoint (gap 0.63pp); s123/s456 use seed-matched converged oracles ($\leq$0.03pp gap).}, while magnitude-trained BCE yields $2.4\times$ worse PPL (Table~\ref{tab:main}).
KL-discovered mask patterns thus \emph{transfer} to simpler training, enabling a practical oracle-then-distill deployment pipeline: extract KL-derived oracles once, then train simpler BCE predictors per model.
BCE+Gumbel (32.22) underperforms standard BCE (27.97), indicating that Gumbel relaxation hurts when the target is a fixed binary mask rather than a distributional objective (Appendix~\ref{app:oracle_provenance}).

\subsection{Ablation Studies}

Table~\ref{tab:ablation} isolates the contribution of each design choice at 30\% and 50\% sparsity.

\begin{table}[t]
\centering
\caption{Ablation on LLaMA-3.1-8B (WikiText-2 PPL$\downarrow$; uniform at 30\%, global at 50\%). $\pm$: 3-seed; unmarked: s42. $^\ddagger$STE: 20.8\% actual sparsity at 30\% target; loss effect ($\Delta$PPL=20.4) is $5.1\times$ relaxation effect ($\Delta$PPL=3.98). $^\S$MSE 50\%: 2/3 seeds collapsed to ${\sim}$5\% actual sparsity; s42 only (Appendix~\ref{app:mse_config}).}
\label{tab:ablation}
\resizebox{\columnwidth}{!}{
\begin{tabular}{lccc}
\toprule
Configuration & 30\% PPL$\downarrow$ & 50\% PPL$\downarrow$ & $\Delta_{50}$ \\
\midrule
\multicolumn{4}{l}{\textit{KL-based variants}} \\
Forward KL + Gumbel (Ours) & \textbf{8.81$\pm$0.02} & \textbf{11.79} & --- \\
JSD (symmetric) + Gumbel & 8.90$\pm$0.02 & 12.54 & +0.75 \\
Reverse KL + Gumbel & 9.26$\pm$0.07 & 12.80 & +1.01 \\
Forward KL + STE & ---$^\ddagger$ & 15.77 & +3.98 \\
\midrule
\multicolumn{4}{l}{\textit{Per-layer training}} \\
Per-layer MSE & 18.52$\pm$1.68 & 229.26$^\S$ & +217.47 \\
Per-layer KL ($K{=}1$) & 9.05$\pm$0.02 & 34.85 & +23.06 \\
Per-layer KL ($K{=}4$) & 8.89$\pm$0.01 & 44.23 & +32.44 \\
\midrule
\multicolumn{4}{l}{\textit{BCE-based variants}} \\
BCE (no compensation) & 15.71 & 27.97 & +16.18 \\
BCE + Gumbel (no penalty) & --- & 32.22 & +20.43 \\
BCE + Gumbel + Penalty & --- & 38.02 & +26.23 \\
BCE + STE & --- & 60.54 & +48.75 \\
\bottomrule
\end{tabular}}
\end{table}

\paragraph{Relaxation and Divergence.}
At 50\%, Gumbel-Sigmoid outperforms STE~\cite{bengio2013ste} by 25.2\%; forward KL outperforms JSD (+6.0\%) and reverse KL (+7.9\%), consistent with mode-covering~\cite{gu2024minillm}.
STE's failure to achieve the target sparsity (actual: 20.8\% at 30\% target) itself reflects a limitation of straight-through gradients for differentiable masking: the hard-threshold forward pass creates gradient mismatch that destabilizes Lagrangian enforcement.
At 30\%, forward KL (8.81) and JSD (8.90) are near-identical, both outperforming reverse KL (9.26) and MSE (18.52$\pm$1.68), suggesting the divergence measure matters less than optimizing through the output distribution (\S\ref{sec:perlayer}).
\paragraph{Per-Layer KL: Signal Type vs.\ Coupling Scope.}
\label{sec:perlayer}
To isolate cross-layer coupling, we train per-layer KL predictors masking $K$ random layers per step, computing KL on the full output without cross-layer gradients.
The $K$-curve (Figure~\ref{fig:k_curve_ppl}) shows per-layer KL ($K{=}1$) achieves 9.05$\pm$0.02, within 2.7\% of end-to-end (E2E) training (8.81), while per-layer MSE degrades to 18.52$\pm$1.68.
The signal type gap\footnote{``Signal type'' denotes the optimization target (output distribution vs.\ intermediate activations), which also entails gradient flow and landscape differences not independently controlled.} (9.47 PPL) exceeds the coupling breadth range (2.08 PPL; full $K$-curve in Appendix~\ref{app:k_curve}), suggesting output-distribution alignment contributes more than coupling scope, though confounded by train--inference mismatch (\S\ref{app:k4_downstream}).
A $K{=}4$ control partially deconfounds: PPL gaps close while commonsense gaps persist, suggesting task-dependent coupling (Appendices~\ref{app:k4_downstream}--\ref{app:perlayer_downstream}).

\begin{figure}[t]
\centering
\includegraphics[width=\columnwidth]{k_curve_ppl}
\vspace{-8pt}
\caption{Coupling breadth $K$ vs.\ perplexity (30\%, TEAL uniform, 3-seed mean$\pm$std). $K{=}32$ denotes end-to-end training. $K{=}2$ degradation reflects insufficient per-layer updates (${\sim}62$ vs.\ ${\sim}250$ at $K{=}8$).}
\label{fig:k_curve_ppl}
\vspace{-8pt}
\end{figure}

\paragraph{Compensation Redundancy.}
\label{sec:compensation}
FastForward-style compensation~\cite{gautam2026fastforward} (+67.2M params) provides no benefit (Appendices~\ref{app:comp_convergence}--\ref{app:training_convergence}); SPON bias correction is orthogonal (Appendix~\ref{app:spon}); predictor capacity and temperature show minimal sensitivity (Appendices~\ref{app:db_sensitivity}--\ref{app:tau_sensitivity}).


\subsection{Sparsity Scaling and Regime Analysis}
\label{sec:scaling}

\begin{figure}[t]
\centering
\includegraphics[width=\columnwidth]{fig2_sparsity_sweep.pdf}
\vspace{-8pt}
\caption{PPL scaling with sparsity (global allocation). KL's advantage over TEAL grows with sparsity.}
\label{fig:sparsity_sweep}
\vspace{-8pt}
\end{figure}

\paragraph{Allocation Behavior.}
Table~\ref{tab:phase} (Appendix~\ref{app:alloc_full}) reveals non-monotonic allocation preference: uniform is best at $\leq$30\% ($7.3\times$ smaller MMLU $\sigma$), global at 35--40\%, both collapse at $\geq$45\%.
Global top-$k$ concentrates budget in high-importance middle layers (L14--18, Appendix~\ref{app:retention}) but starves late layers (L31 retains only 28.4\% at 30\%); uniform avoids this pathology when the total budget is sufficient.

\paragraph{Cross-Architecture and Scale.}
\label{sec:cross_arch}
On Mistral-7B-v0.3~\cite{jiang2023mistral}, PPL generalizes ($-$17.3\%, 3-seed; Appendix~\ref{app:cross_arch}) with directional commonsense improvements; on Qwen-2.5-14B~\cite{qwen2024qwen25}, PPL generalizes but MMLU and GSM8K decline (Appendix~\ref{app:14b}).
The knowledge regression traces to allocation-induced layer starvation (uniform reduces MMLU variance $7.3\times$; Appendix~\ref{app:constraint}) and the objective's bias toward high-probability over rare factual tokens~\cite{geva2023dissecting}.
The pipeline adds 6.4\% latency overhead absent sparse kernels (Appendix~\ref{app:latency}).

